\documentclass[11pt]{article}

\usepackage[margin=25mm]{geometry}

\usepackage{amsmath,amssymb,fontspec,unicode-math}

\IfFontExistsTF{DejaVu Serif}{\setmainfont{DejaVu Serif}}{\setmainfont{Latin Modern Roman}}

\IfFontExistsTF{Latin Modern Math}{\setmathfont{Latin Modern Math}}{\setmathfont{DejaVu Math TeX Gyre}}

\usepackage{xurl,hyperref}

\setlength{\parindent}{0pt}\setlength{\parskip}{6pt}

\begin{document}

\title{Enhancing Infrared Vision: Progressive Prompt Fusion Network and Benchmark}

\author{ELVIA preliminary AI analysis}\date{21 October 2026}\maketitle

Preliminary assessment: Minor.\par

Independent scientific review remains pending. Author adoption has not been requested.

Singleton prompt dimensions make the released multi-device training path discard image shards and can break validation. The stated low-contrast formula also differs from its pinned generator’s brightness offset. Historical experiment correspondence remains unresolved.

\section*{Review scope}

Claim-focused reading of the original paper and relevant author-released passages; direct mathematical or static-source reasoning. No author program, earlier verifier, model or benchmark was executed during this check.

\section*{Problem 1: A singleton prompt truncates parallel image batches}

Source: train.py, lines 146–156,188; PyTorch 1.12 scatter\_map

Assessment: supported; impact: minor

\subsection*{Source claim (paraphrase)}

Training passes images shaped \((4,1,256,256)\), a degradation prompt \((1,128,128)\), and scenario prompts \((4,64)\) into DataParallel.

\subsection*{Analysis}

Dimension-zero scattering yields four image/scenario chunks but one degradation chunk. The recursive zip stops at one input, leaving a one-image prediction. For scalar predictor \(f_w(x)=wx\), \(w=1\), four constant images/targets \(0.1,0.2,0.7,0.8\) should have zero loss. Broadcasting the surviving first prediction gives mean absolute loss \(0.35\) and gradient \(-0.075\). Validation passes a scenario vector \((64,)\), which is split into length \(16\) on four destinations, incompatible with the released \(64\)-input linear layer. Single-device operation can avoid this scattering defect.

\subsection*{Finding and scope}

The released multi-device call layout fails its per-example contract. This is a static source/shape finding; no claim of four-physical-GPU execution or false published scores is made.

\subsection*{Proposed correction}

Before scattering, give every tensor a matching batch axis: degradation \((B,128,128)\), scenario \((B,64)\). Assert equal prediction/target shapes. This restores the identified correspondence, not end-to-end trained performance.

A proposed correction is retained where stated; independent verification and author adoption remain unestablished.

\section*{Problem 2: The corruption’s brightness offset differs from the printed formula}

Source: Appendix A.1.1 Eq.(8), PDF p.22; ranges, PDF p.23; dataset.py/create\_corrupt.py and Albumentations 1.4.1

Assessment: supported; impact: minor

\subsection*{Source claim (paraphrase)}

The paper specifies \(C(I)=\alpha(I+\beta M)\); the pinned brightness-by-max transform implements \(\alpha I+\beta M\).

\subsection*{Analysis}

Use valid uint8 intensities \(0,100,200,255\), \(M=255\), and interior hard-range values \(\alpha=0.5,\beta=0.3\). The printed real outputs are \[(38.25,88.25,138.25,165.75)\]; the dependency gives \((76.5,126.5,176.5,204)\). No value saturates, and the real offset difference is \((1-\alpha)\beta M=38.25\). Quantization cannot explain it. Renaming brightness would require correlated parameter changes; at these values the implied paper brightness \(0.6\) exceeds the hard range.

\subsection*{Finding and scope}

The release and printed parameters define different corrupted inputs. Which produced the historical benchmark images is unestablished.

\subsection*{Proposed correction}

Either document the multiply-then-add generator and ranges, or implement the printed correlated offset. Matching one formula locally does not establish improved restoration or recover historical data.

A proposed correction is retained where stated; independent verification and author adoption remain unestablished.

\section*{Limits}

Preliminary AI assessment; independent human scientific adjudication remains pending.\par

Analytic witnesses and static release checks do not remeasure the published benchmark results.\par

\section*{Source identity}

NeurIPS 2025 official proceedings PDF (retained local copy)

\url{https://proceedings.neurips.cc/paper_files/paper/2025/file/8ba8b021e0f7f6dd6262227214e1517f-Paper-Conference.pdf}

Reviewed PDF SHA-256: \nolinkurl{d912d6fc6880d3146fd22933e66c82259ae143b3bf2395f05d99cad128c92af7}

The SHA-256 identifies the exact retained PDF inspected in this rapid source/argument check. Original supplied report bytes are retained locally. The original PDF retrieval timestamp is taken from matching retained metadata.

\end{document}
